\section{Discussion}
\label{sec:discussion}

\para{Person space is low-rank and shared.}
Our central finding is that LLM person space has approximately \effrank{} effective dimensions---roughly one-third of the full preference space---and that this structure is nearly identical across model sizes ($r = 0.96$).
This suggests that the implicit model of ``which preferences go together'' is not an artifact of a particular model's capacity or training details, but a robust feature of the training data distribution itself.
The high cross-model consistency ($r > 0.90$ for all top-5 components) is particularly striking given that \nano and \mini differ substantially in parameter count and capability.

\para{Relationship to activation-level findings.}
\citet{chen2025persona} and \citet{bhandari2026personality} find entangled, non-orthogonal personality directions in activation space, with a dominant ``social desirability'' axis.
Our Component~1 (outgoing intellectual vs.\ homebody consumer) may reflect a behavioral analog of this axis.
The partial interpretability of our components---they capture real structure but do not map cleanly onto Big Five---mirrors the geometric findings of \citet{bhandari2026personality}, who show that behavioral independence cannot be achieved even after geometric orthogonalization.
A direct comparison between behavioral and activation-level dimensionality remains an important open question.

\para{Implications for personalization.}
If person space is truly \effrank-dimensional, then a user's full preference profile can in principle be inferred from a small number of observations.
This connects to the few-shot adaptation results of \citet{bose2025lore}, who show that 3--4 preference comparisons suffice to locate a user in reward space.
Our findings suggest that similar efficiency should hold for behavioral personalization: the low-rank structure means that preferences are highly predictable from a compact set of dimensions.

\para{Implications for alignment and safety.}
The shared person-space structure implies that LLMs have strong priors about what constitutes a ``coherent person.''
Profiles that violate this structure---those far from the low-dimensional manifold---produce inconsistent behavior, which could be exploited or could cause unexpected failures in applications relying on consistent persona maintenance.
Understanding these constraints helps predict where persona-based safety measures (\eg ``act as a helpful, harmless assistant'') may break down.

\subsection{Limitations}
\label{sec:limitations}

\para{Self-evaluation bias.}
Both congruence scoring and coherence rating use models from the same family (GPT-4.1) as the persona-adopting models.
This may inflate scores and reduce measured variance compared to cross-family evaluation.
Future work should use independent judges (\eg human raters or models from different families) to validate our measurements.

\para{Limited model diversity.}
We test only two models from the same provider (OpenAI).
While the high cross-model correlation is encouraging, testing models from different families (Claude, Gemini, open-source models) is needed to determine whether the $r = 0.96$ structural similarity reflects shared training data, shared architecture, or a more universal feature of language modeling.

\para{Sparse profile-preference matrix.}
With \nprofiles{} profiles of 8 items each drawn from \nprefs{} possibilities, the profile-preference matrix is 87\% zeros.
Our analysis focuses on the denser co-occurrence matrix, but richer coverage (more profiles or more items per profile) would improve estimates of pairwise interactions.

\para{Preference item validity.}
Our \nprefs{} items, while diverse, are not drawn from validated psychological instruments.
Different item sets could yield different effective dimensionalities.
Future work should compare our approach with established instruments (\eg items derived from Big Five facets or the HEXACO model) to calibrate the discovered dimensions.

\para{Scoring noise.}
Using an LLM to score congruence introduces systematic biases.
Temperature 0.7 in persona responses adds stochastic variation.
Multiple independent runs per profile and human evaluation of a subset would strengthen reliability.
